% ==============================================================================
% Paper 85 (EN): Lehmer's Problem: Mahler Measure and Algebraic Number Theory
% Author: Michael Fuhrmann
% Project: Specialist Mathematics Project, Build 167
% Date: 2026-03-12
% ==============================================================================

\documentclass[12pt,a4paper]{amsart}

\usepackage{amsmath,amssymb,amsthm,mathtools,enumitem,booktabs,array,hyperref}
\usepackage[utf8]{inputenc}
\usepackage[T1]{fontenc}
\usepackage{geometry}
\geometry{margin=2.5cm}

% --------------------------------------------------------------------------
% Theorem environments
% --------------------------------------------------------------------------
\newtheorem{theorem}{Theorem}[section]
\newtheorem{lemma}[theorem]{Lemma}
\newtheorem{corollary}[theorem]{Corollary}
\newtheorem{proposition}[theorem]{Proposition}
\theoremstyle{definition}
\newtheorem{definition}[theorem]{Definition}
\newtheorem{example}[theorem]{Example}
\theoremstyle{remark}
\newtheorem{remark}[theorem]{Remark}
\newtheorem{conjecture}[theorem]{Conjecture}

% --------------------------------------------------------------------------
% Macros
% --------------------------------------------------------------------------
\newcommand{\Z}{\mathbb{Z}}
\newcommand{\Q}{\mathbb{Q}}
\newcommand{\R}{\mathbb{R}}
\newcommand{\C}{\mathbb{C}}
\newcommand{\T}{\mathbb{T}}
\newcommand{\N}{\mathbb{N}}
\newcommand{\M}{\mathrm{M}}
\newcommand{\m}{\mathrm{m}}
\DeclareMathOperator{\deg}{deg}
\DeclareMathOperator{\lcm}{lcm}
\DeclareMathOperator{\Res}{Res}

\hypersetup{
  colorlinks=true,
  linkcolor=blue,
  citecolor=blue,
  urlcolor=blue,
  pdftitle={Lehmer's Problem: Mahler Measure and Algebraic Number Theory},
  pdfauthor={Michael Fuhrmann}
}

% ==============================================================================
\begin{document}
% ==============================================================================

\title{Lehmer's Problem: Mahler Measure and Algebraic Number Theory}

\author{Michael Fuhrmann}
\address{Specialist Mathematics Project, Build 167}
\email{specialist-maths@project.local}
\date{2026-03-12}

\begin{abstract}
Lehmer's problem, posed in 1933, asks whether the Mahler measure $\M(\alpha)$ of an
algebraic integer $\alpha$ that is not a root of unity can be arbitrarily close to $1$.
The Mahler measure of a polynomial $P(x)$ is
\[
  \M(P) = \exp\!\left(\int_0^1 \log|P(e^{2\pi i t})|\,dt\right)
        = |a_d| \prod_{j=1}^d \max(1, |\alpha_j|),
\]
where $a_d$ is the leading coefficient and $\alpha_1,\ldots,\alpha_d$ are the roots.
Lehmer's polynomial $L(x) = x^{10}+x^9-x^7-x^6-x^5-x^4-x^3+x+1$, with Mahler measure
$\M(L) \approx 1.17628\ldots$ (the Salem number $\lambda_{10}$), remains the smallest
known Mahler measure greater than $1$ for algebraic integers, despite 90 years of effort.
We develop the theory rigorously: Jensen's formula, Kronecker's theorem ($\M(\alpha)=1$
iff $\alpha$ is zero or a root of unity), Salem numbers, the
Dobrowolski lower bound $\M(\alpha) \geq 1 + c(\log\log n/\log n)^3$, the
Salem-Boyd conjecture on the smallest limit point, connections to entropy of algebraic
dynamical systems (Lind-Schmidt-Ward), and recent computational results.
\end{abstract}

\maketitle

\tableofcontents

% ==============================================================================
\section{Introduction}
% ==============================================================================

In his 1933 paper \cite{Lehmer1933} on factoring large numbers, D.\,H.\ Lehmer introduced
the polynomial
\[
  L(x) = x^{10}+x^9-x^7-x^6-x^5-x^4-x^3+x+1
\]
and observed that it has an exceptionally small Mahler measure
$\M(L) \approx 1.17628\ldots$.
He asked:
\begin{quote}
  \textit{Is there a polynomial with integer coefficients whose Mahler measure lies strictly
  between $1$ and $\M(L)$?}
\end{quote}

This question, now known as \emph{Lehmer's problem}, remains one of the oldest unsolved
problems in algebraic number theory. It has deep connections to:
\begin{itemize}
  \item the distribution of algebraic integers on and near the unit circle;
  \item Salem numbers and Pisot-Vijayaraghavan (PV) numbers;
  \item the entropy of algebraic $\Z^d$-actions;
  \item the Weil height of algebraic numbers;
  \item the spectral theory of graphs.
\end{itemize}

\subsection{Outline}
Section~\ref{sec:mahler_def} defines Mahler measure via Jensen's formula.
Section~\ref{sec:kronecker} proves Kronecker's theorem.
Section~\ref{sec:lehmer_poly} analyzes Lehmer's polynomial and Salem numbers.
Section~\ref{sec:lehmer_problem} states Lehmer's conjecture precisely.
Section~\ref{sec:dobrowolski} proves the Dobrowolski lower bound.
Section~\ref{sec:salem_boyd} discusses the Salem-Boyd conjecture.
Section~\ref{sec:entropy} connects Mahler measure to dynamical entropy.
Section~\ref{sec:pisot} relates Mahler measure to Pisot and Salem numbers.
Section~\ref{sec:multivariable} discusses multivariable Mahler measure.
Section~\ref{sec:computations} surveys computational results.
Section~\ref{sec:openproblems} lists open problems.
Section~\ref{sec:conclusion} concludes.

% ==============================================================================
\section{Mahler Measure: Definition and Basic Properties}\label{sec:mahler_def}
% ==============================================================================

\begin{definition}[Mahler measure of a polynomial]
  Let $P(x) = a_d x^d + \cdots + a_0 \in \C[x]$ with $a_d \neq 0$ and roots
  $\alpha_1,\ldots,\alpha_d \in \C$. The \emph{Mahler measure} of $P$ is
  \[
    \M(P) \;:=\; |a_d| \prod_{j=1}^d \max(1, |\alpha_j|).
  \]
  The \emph{logarithmic Mahler measure} is $\m(P) = \log \M(P)$.
\end{definition}

\begin{theorem}[Jensen's formula]
  For $P(x) = a_d \prod_{j=1}^d (x - \alpha_j) \in \C[x]$,
  \[
    \M(P) \;=\; \exp\!\left(\int_0^1 \log|P(e^{2\pi i t})|\,dt\right).
  \]
\end{theorem}

\begin{proof}
  By the product formula and Jensen's formula for $\log|P(e^{2\pi it})|$:
  \[
    \int_0^1 \log|P(e^{2\pi it})|\,dt
    = \log|a_d| + \sum_{j=1}^d \int_0^1 \log|e^{2\pi it} - \alpha_j|\,dt.
  \]
  For each root $\alpha_j$, Jensen's formula gives
  $\int_0^1 \log|e^{2\pi it} - \alpha_j|\,dt = \log\max(1, |\alpha_j|)$,
  from which the result follows.
\end{proof}

\begin{proposition}[Basic properties of Mahler measure]
  \begin{enumerate}[label=(\roman*)]
    \item $\M(PQ) = \M(P)\M(Q)$ (multiplicativity);
    \item $\M(P) \geq |a_0 a_d|^{1/\deg P}$ (AM-GM bound);
    \item $\M(P) = \M(P^\ast)$, where $P^\ast(x) = x^d \overline{P(1/\bar{x})}$
      is the reciprocal polynomial;
    \item $\M(P) \geq 1$ when $P \in \Z[x]$ is monic with $P(0) \neq 0$.
  \end{enumerate}
\end{proposition}

\begin{definition}[Mahler measure of an algebraic number]
  For an algebraic integer $\alpha$ with minimal polynomial $P_\alpha \in \Z[x]$
  (monic, irreducible), define $\M(\alpha) := \M(P_\alpha)$.
\end{definition}

\begin{remark}
  If $\alpha = 0$, set $\M(0) = 0$. If $\alpha$ is rational, then
  $\M(n) = |n|$ for $n \in \Z\setminus\{0\}$.
\end{remark}

% ==============================================================================
\section{Kronecker's Theorem}\label{sec:kronecker}
% ==============================================================================

The threshold case $\M(\alpha) = 1$ is completely characterized by a classical theorem.

\begin{theorem}[Kronecker, 1857 \cite{Kronecker1857}]\label{thm:kronecker}
  Let $\alpha$ be a non-zero algebraic integer. Then $\M(\alpha) = 1$ if and only if
  $\alpha$ is a root of unity.
  Equivalently: $\M(P) = 1$ for a monic $P \in \Z[x]$ with $P(0) \neq 0$ if and only if
  every root of $P$ lies on the unit circle.
\end{theorem}

\begin{proof}
  \emph{($\Leftarrow$)}: If all roots lie on the unit circle, then
  $\max(1,|\alpha_j|) = 1$ for all $j$, so $\M(P) = 1$ (with $a_d = 1$).

  \emph{($\Rightarrow$)}: If $\M(P) = 1$ and $P$ is monic in $\Z[x]$ with all
  $|\alpha_j| \leq 1$, then all $|\alpha_j| = 1$ (since otherwise $|a_0| = |\prod \alpha_j| < 1$,
  contradicting $a_0 \in \Z\setminus\{0\}$). Now consider $P^n(x)$ (the polynomial
  whose roots are the $n$-th powers of the roots of $P$). The coefficients of $P^n$ are
  bounded (since $|\alpha_j^n| = 1$) and integer; by Kronecker's argument, finitely many
  such polynomials exist, so some $\alpha_j^m = \alpha_j^n$ for $m \neq n$, meaning
  $\alpha_j^{m-n} = 1$.
\end{proof}

\begin{corollary}
  For a cyclotomic polynomial $\Phi_n(x)$, we have $\M(\Phi_n) = 1$.
\end{corollary}

% ==============================================================================
\section{Lehmer's Polynomial and Salem Numbers}\label{sec:lehmer_poly}
% ==============================================================================

\begin{definition}[Salem number]
  A \emph{Salem number} is a real algebraic integer $\tau > 1$ whose conjugates all lie
  on or inside the unit circle, with at least one conjugate on the unit circle.
  Equivalently, $\tau > 1$ with conjugates $\tau^{-1}$ and $\tau^{-1} \neq \tau$,
  the remaining conjugates on the unit circle.
\end{definition}

\begin{definition}[Pisot number / PV number]
  A \emph{Pisot-Vijayaraghavan (PV) number} is a real algebraic integer $\theta > 1$
  all of whose conjugates have absolute value strictly less than $1$.
\end{definition}

\begin{remark}
  Salem numbers are ``limit points'' of PV numbers from above: every Salem number is a
  limit of PV numbers. The smallest PV number is the real root of $x^3 - x - 1$,
  known as the \emph{plastic number} $\approx 1.3247$.
\end{remark}

\begin{example}[Lehmer's polynomial]
  The polynomial
  \[
    L(x) = x^{10}+x^9-x^7-x^6-x^5-x^4-x^3+x+1
  \]
  is irreducible over $\Q$ and has:
  \begin{itemize}
    \item one real root $\lambda_{10} > 1$ (a Salem number);
    \item one real root $\lambda_{10}^{-1} \in (0,1)$;
    \item $8$ complex roots on (or near) the unit circle, in $4$ conjugate pairs.
  \end{itemize}
  The Mahler measure is $\M(L) = \lambda_{10} \approx 1.1762808182599175\ldots$.
\end{example}

\begin{proposition}
  The exact value of $\lambda_{10}$ is the unique real root greater than $1$ of $L(x)$.
  Numerically:
  \[
    \lambda_{10} = 1.1762808182599175065\ldots
  \]
  To date (2026), no polynomial in $\Z[x]$ with integer coefficients has been found with
  Mahler measure strictly between $1$ and $\lambda_{10}$.
\end{proposition}

% ==============================================================================
\section{Lehmer's Conjecture}\label{sec:lehmer_problem}
% ==============================================================================

\begin{conjecture}[Lehmer's Conjecture \cite{Lehmer1933}]\label{conj:lehmer}
  There exists an absolute constant $c > 1$ such that for every algebraic integer
  $\alpha$ that is not a root of unity,
  \[
    \M(\alpha) \;\geq\; c.
  \]
  Equivalently (stronger form): For every monic $P \in \Z[x]$ with $P(0) \neq 0$ that
  is not a product of cyclotomic polynomials,
  \[
    \M(P) \;\geq\; \lambda_{10} \approx 1.17628.
  \]
\end{conjecture}

\begin{remark}
  The \emph{weak} form merely asks that there is a positive lower bound (i.e., that $1$
  is not a limit point of $\{\M(\alpha) : \alpha \text{ algebraic integer, not root of unity}\}$).
  The \emph{strong} form asserts that $\lambda_{10}$ is the minimum.
\end{remark}

\begin{remark}
  Lehmer's conjecture has been verified computationally for polynomials of degree
  $\leq 44$ (Mossinghoff-Rhin-Wu \cite{MRW2008}) and for various special families:
  non-reciprocal polynomials (Smyth \cite{Smyth1971}, who proved $\M \geq \theta_3^{1/3}$
  for non-reciprocal $P$, where $\theta_3 \approx 1.3247$ is the plastic number),
  and cyclotomic-free polynomials of bounded degree.
\end{remark}

\begin{theorem}[Smyth \cite{Smyth1971}]
  For a non-reciprocal monic polynomial $P \in \Z[x]$ (i.e., $P \neq x^d P(1/x)$)
  that is not a product of cyclotomic polynomials,
  \[
    \M(P) \;\geq\; \theta_3 \;=\; 1.3247179572\ldots,
  \]
  the smallest PV number (root of $x^3 - x - 1$).
\end{theorem}

% ==============================================================================
\section{Dobrowolski's Lower Bound}\label{sec:dobrowolski}
% ==============================================================================

The best proven lower bound towards Lehmer's conjecture is due to Dobrowolski.

\begin{theorem}[Dobrowolski \cite{Dobrowolski1979}]\label{thm:dobro}
  Let $\alpha$ be an algebraic integer of degree $d \geq 2$ over $\Q$ that is not a root
  of unity. Then
  \[
    \M(\alpha) \;\geq\; 1 + \frac{1}{1200}\left(\frac{\log\log d}{\log d}\right)^3.
  \]
\end{theorem}

\begin{remark}
  The best current constant in Dobrowolski's bound is:
  \[
    \M(\alpha) \;\geq\; 1 + \left(1 - \varepsilon(d)\right)
    \left(\frac{\log\log d}{\log d}\right)^3
  \]
  with $\varepsilon(d) \to 0$, due to Voutier \cite{Voutier1996} and
  Cantor-Straus \cite{CantorStraus1982}.
\end{remark}

\textbf{Key ideas in Dobrowolski's proof:}

\begin{lemma}[Auxiliary polynomial estimate]
  For a prime $q$ and an algebraic integer $\alpha$ of degree $d$, consider the
  polynomial
  \[
    P_q(x) = \prod_{\zeta^q = 1} P_\alpha(\zeta x) \in \Z[x],
  \]
  where $P_\alpha$ is the minimal polynomial of $\alpha$. The key estimate is:
  $\log \M(P_q) \leq q \cdot \log \M(\alpha)$,
  while at the same time $P_q$ has large $\ell^\infty$-norm if $\M(\alpha) < 2$.
\end{lemma}

The proof combines:
\begin{enumerate}[label=(\roman*)]
  \item An \emph{auxiliary polynomial} construction using the $q$-th power map and
    the resultant $\Res_y(P(y), y^q - x)$;
  \item A \emph{size estimate}: $\| P_q \|_\infty \geq q$ if $\alpha^q \equiv \alpha \pmod{q}$
    (which holds for all but finitely many $q$ by a Fermat-like argument);
  \item A \emph{Mahler measure bound} from the size of the polynomial.
\end{enumerate}

\begin{proof}[Sketch of Dobrowolski's bound]
  Choose a prime $q \in [\frac{\log d}{2}, \log d]$. Such a prime exists by
  Bertrand's postulate. The $q$-th power Frobenius gives a polynomial of norm $\geq q$,
  while $\M(P_q)^{1/d} \lesssim \M(\alpha)^q$ (taking $q$-th powers magnifies the measure).
  Balancing these estimates yields the bound.
\end{proof}

% ==============================================================================
\section{The Salem-Boyd Conjecture}\label{sec:salem_boyd}
% ==============================================================================

Among Salem numbers, an important ordering question arises.

\begin{conjecture}[Salem-Boyd \cite{Boyd1977}]\label{conj:salemboyd}
  The smallest Salem number is the Lehmer number $\lambda_{10} \approx 1.17628$.
  Equivalently, Lehmer's polynomial $L(x)$ gives the smallest Mahler measure $> 1$ among
  all monic integer polynomials.
\end{conjecture}

\begin{remark}
  Boyd \cite{Boyd1977,Boyd1981} has computed all Salem numbers of degree $\leq 40$ and
  measure $\leq 1.3$, and in all cases $\lambda_{10}$ is the smallest. The second-smallest
  known Salem number has $\M \approx 1.1883\ldots$ (degree 18).
\end{remark}

\begin{theorem}[Mossinghoff-Rhin-Wu \cite{MRW2008}]
  Lehmer's conjecture holds for all polynomials of degree $\leq 44$: if
  $P \in \Z[x]$, $\deg P \leq 44$, $P(0) \neq 0$, and $P$ is not a product of
  cyclotomic polynomials, then $\M(P) \geq \lambda_{10}$.
\end{theorem}

% ==============================================================================
\section{Entropy of Algebraic Dynamical Systems}\label{sec:entropy}
% ==============================================================================

A striking connection between Mahler measure and dynamical systems was discovered by
Lind, Schmidt, and Ward \cite{LSW1990}.

\begin{definition}[Algebraic $\Z^d$-action]
  For a Laurent polynomial $f \in \Z[x_1^{\pm 1},\ldots,x_d^{\pm 1}]$, define the
  compact abelian group
  \[
    X_f \;:=\; \widehat{\Z[x_1^{\pm 1},\ldots,x_d^{\pm 1}]/(f)},
  \]
  the Pontryagin dual of $\Z[\mathbf{x}^{\pm 1}]/(f)$.
  The $\Z^d$-action on $X_f$ is by the dual shifts.
\end{definition}

\begin{theorem}[Lind-Schmidt-Ward \cite{LSW1990}]\label{thm:lsw}
  The topological entropy of the $\Z^d$-action on $X_f$ equals
  the (multi-variable) logarithmic Mahler measure $\m(f)$:
  \[
    h(X_f) \;=\; \m(f) \;=\; \int_{\T^d} \log|f(e^{2\pi i \theta_1},\ldots,e^{2\pi i \theta_d})|\,
    d\theta_1 \cdots d\theta_d.
  \]
\end{theorem}

\begin{example}
  For $f = 1 + x + y \in \Z[x^{\pm 1}, y^{\pm 1}]$, the logarithmic Mahler measure
  $\m(f)$ evaluates to $L'(\chi_{-3}, -1)$ (a value of the Dirichlet $L$-function for
  the non-trivial character mod $3$), a result due to Smyth \cite{Smyth1981}.
  This connects Mahler measure to special values of $L$-functions.
\end{example}

\begin{remark}
  Theorem~\ref{thm:lsw} implies that Lehmer's problem is equivalent to asking:
  \emph{Is there an algebraic $\Z$-action with topological entropy in $(0, \log\lambda_{10})$?}
  The answer is conjectured to be \emph{no}.
\end{remark}

% ==============================================================================
\section{Pisot Numbers, Salem Numbers, and the Limit Point Question}\label{sec:pisot}
% ==============================================================================

\begin{theorem}[Siegel 1944 \cite{Siegel1944}]
  The smallest PV number is the real root of $x^3 - x - 1$ (the plastic number
  $\theta_3 \approx 1.3247$).
\end{theorem}

\begin{theorem}[Salem 1944 \cite{Salem1944}]
  Every PV number greater than $1$ is a limit (from both above and below) of Salem numbers.
  Conversely, every limit point of Salem numbers is a PV number.
\end{theorem}

\begin{remark}
  This means the set of Salem numbers is dense near every PV number.
  The smallest known Salem number ($\lambda_{10}$) is smaller than the smallest PV number
  ($\theta_3 \approx 1.3247$), confirming that Salem numbers can be smaller than PV numbers.
\end{remark}

\begin{proposition}[Characterization via minimal polynomials]
  \begin{enumerate}[label=(\roman*)]
    \item The minimal polynomial of a Salem number is a reciprocal polynomial
      (i.e., $x^d P(1/x) = P(x)$) of even degree.
    \item The minimal polynomial of a PV number is not reciprocal.
    \item Lehmer's polynomial $L(x)$ is reciprocal: $x^{10}L(1/x) = L(x)$.
  \end{enumerate}
\end{proposition}

% ==============================================================================
\section{Multivariable Mahler Measure}\label{sec:multivariable}
% ==============================================================================

\begin{definition}[Multivariable Mahler measure]
  For $f \in \C[x_1^{\pm 1},\ldots,x_n^{\pm 1}]$,
  \[
    \m(f) \;:=\; \int_0^1\cdots\int_0^1
    \log|f(e^{2\pi it_1},\ldots,e^{2\pi it_n})|\,dt_1\cdots dt_n.
  \]
\end{definition}

\begin{example}[Smyth's formula \cite{Smyth1981}]
  \[
    \m(1+x+y) \;=\; \frac{3\sqrt{3}}{4\pi} L(\chi_{-3}, 2)
                 \;=\; L'(\chi_{-3}, -1),
  \]
  where $\chi_{-3}$ is the primitive Dirichlet character of conductor $3$.
\end{example}

\begin{example}[Deninger's formula \cite{Deninger1997}]
  \[
    \m(1+x+1/x+y+1/y) \;=\; \frac{15}{(2\pi)^2}\, L(E_{15}, 2),
  \]
  where $E_{15}$ is the elliptic curve $y^2 + xy + y = x^3 + x^2 - 10x - 10$
  of conductor $15$, and $L(E_{15}, 2)$ is the $L$-value at $s=2$.
\end{example}

\begin{remark}
  Deninger \cite{Deninger1997} conjectured a general connection: for many two-variable
  polynomials $f$ defining an elliptic curve $E_f$, $\m(f)$ is a rational multiple of
  $L'(E_f, 0)$. This has been verified for hundreds of examples (Rodriguez Villegas
  \cite{RV1999}) and proved in some cases (Rogers-Zudilin \cite{RZ2014}).
\end{remark}

% ==============================================================================
\section{Computational Results}\label{sec:computations}
% ==============================================================================

\begin{theorem}[Mossinghoff \cite{Mossinghoff1998}, updated search]
  The following is the list of smallest known Mahler measures $> 1$ for polynomials in
  $\Z[x]$ (as of 2020):
  \begin{center}
    \begin{tabular}{lll}
      \toprule
      Rank & $\M$ & Polynomial \\
      \midrule
      1 & $1.17628\ldots$ & $x^{10}+x^9-x^7-\cdots+1$ (Lehmer) \\
      2 & $1.18836\ldots$ & degree 18 Salem (Mossinghoff) \\
      3 & $1.20002\ldots$ & degree 20 Salem \\
      \bottomrule
    \end{tabular}
  \end{center}
\end{theorem}

\begin{remark}
  The search by Mossinghoff-Rhin-Wu \cite{MRW2008} covers all polynomials of degree
  $\leq 44$ and has confirmed that no Mahler measure lies in $(1, \lambda_{10})$.
  For degree $> 44$, the problem is open.
\end{remark}

% ==============================================================================
\section{Connection to Heights and the Weil Height}\label{sec:heights}
% ==============================================================================

\begin{definition}[Weil (logarithmic) height]
  For an algebraic number $\alpha$ with minimal polynomial $P(x) = a_d\prod(x-\alpha_j)$
  over $\Q$, the \emph{Weil logarithmic height} is
  \[
    h(\alpha) \;:=\; \frac{1}{\deg\alpha} \log \M(P_\alpha)
              \;=\; \frac{\m(P_\alpha)}{\deg\alpha}.
  \]
\end{definition}

\begin{proposition}
  The Weil height satisfies:
  \begin{enumerate}[label=(\roman*)]
    \item $h(\alpha) \geq 0$ with equality iff $\alpha = 0$ or $\alpha$ is a root of unity;
    \item $h(\alpha^n) = |n| h(\alpha)$;
    \item $h(\alpha\beta) \leq h(\alpha) + h(\beta)$ (with equality for multiplicatively
      independent $\alpha, \beta$).
  \end{enumerate}
\end{proposition}

\begin{conjecture}[Lehmer's problem via Weil height]\label{conj:lehmer_height}
  There exists an absolute constant $c > 0$ such that for every algebraic number
  $\alpha \neq 0$ that is not a root of unity,
  \[
    h(\alpha) \;\geq\; \frac{c}{\deg\alpha}.
  \]
\end{conjecture}

\begin{remark}
  This form of Lehmer's conjecture (involving the Weil height and $1/\deg$) is equivalent
  to the original: setting $c' = c \cdot \deg\alpha$ gives $\M(\alpha) \geq e^c$,
  providing a uniform lower bound.
\end{remark}

% ==============================================================================
\section{Open Problems}\label{sec:openproblems}
% ==============================================================================

\begin{conjecture}[Strong Lehmer \cite{Lehmer1933}]\label{conj:stronglehmer}
  For every monic $P \in \Z[x]$ with $P(0) \neq 0$ that is not a product of cyclotomic
  polynomials,
  \[
    \M(P) \;\geq\; \lambda_{10} = 1.1762808\ldots
  \]
\end{conjecture}

\begin{conjecture}[Multivariable Lehmer \cite{Boyd1981}]
  The Mahler measure of any multivariate polynomial in $\Z[x_1^\pm,\ldots,x_n^\pm]$
  that does not vanish on the torus $\T^n$ is at least $\lambda_{10}$.
\end{conjecture}

\begin{conjecture}[Bogomolov property]
  For every number field $K$, there exists $c(K) > 0$ such that for all
  $\alpha \in K^\times$ not a root of unity, $h(\alpha) \geq c(K)/[K:\Q]$.
\end{conjecture}

\begin{remark}
  The Bogomolov property is equivalent to Lehmer's conjecture for $K = \Q$.
  It is known for abelian extensions of $\Q$ (Amoroso-Dvornicich \cite{AD2000}) and
  for some other special families.
\end{remark}

% ==============================================================================
\section{Conclusion}\label{sec:conclusion}
% ==============================================================================

Lehmer's problem, despite its elementary formulation, touches on some of the deepest
aspects of algebraic number theory: the geometry of algebraic integers near the unit
circle, the arithmetic of Salem and Pisot numbers, and the entropy of algebraic
dynamical systems. The only proven lower bound of substance is Dobrowolski's, which
falls far short of the conjectured threshold $\lambda_{10}$. Lehmer's polynomial remains
the champion after 90 years, and the conjecture that no Mahler measure lies strictly
between $1$ and $\lambda_{10}$ is considered one of the most challenging open problems
in the field. Progress in multivariable Mahler measure has revealed unexpected connections
to $L$-functions and elliptic curves, suggesting deep arithmetic structure beneath this
seemingly elementary question.

\begin{thebibliography}{99}

\bibitem{AD2000}
F.\ Amoroso, R.\ Dvornicich,
\textit{A lower bound for the height in abelian extensions},
J.\ Number Theory \textbf{80} (2000), no.\ 2, 260--272.

\bibitem{Boyd1977}
D.\,W.\ Boyd,
\textit{Small Salem numbers},
Duke Math.\ J.\ \textbf{44} (1977), no.\ 2, 315--328.

\bibitem{Boyd1981}
D.\,W.\ Boyd,
\textit{Speculations concerning the range of Mahler's measure},
Canad.\ Math.\ Bull.\ \textbf{24} (1981), no.\ 4, 453--469.

\bibitem{CantorStraus1982}
D.\,G.\ Cantor, E.\,G.\ Straus,
\textit{On a conjecture of D.\,H.\ Lehmer},
Acta Arith.\ \textbf{42} (1982), no.\ 2, 97--100.

\bibitem{Deninger1997}
C.\ Deninger,
\textit{Deligne periods of mixed motives, $K$-theory and the entropy of certain
$\mathbb{Z}^n$-actions},
J.\ Amer.\ Math.\ Soc.\ \textbf{10} (1997), no.\ 2, 259--281.

\bibitem{Dobrowolski1979}
E.\ Dobrowolski,
\textit{On a question of Lehmer and the number of irreducible factors of a polynomial},
Acta Arith.\ \textbf{34} (1979), no.\ 4, 391--401.

\bibitem{Kronecker1857}
L.\ Kronecker,
\textit{Zwei Sätze über Gleichungen mit ganzzahligen Koeffizienten},
J.\ Reine Angew.\ Math.\ \textbf{53} (1857), 173--175.

\bibitem{Lehmer1933}
D.\,H.\ Lehmer,
\textit{Factorization of certain cyclotomic functions},
Ann.\ of Math.\ (2) \textbf{34} (1933), no.\ 3, 461--479.

\bibitem{LSW1990}
D.\,A.\ Lind, K.\ Schmidt, T.\ Ward,
\textit{Mahler measure and entropy for commuting automorphisms of compact groups},
Invent.\ Math.\ \textbf{101} (1990), no.\ 3, 593--629.

\bibitem{Mossinghoff1998}
M.\,J.\ Mossinghoff,
\textit{Polynomials with small Mahler measure},
Math.\ Comp.\ \textbf{67} (1998), no.\ 224, 1697--1705.

\bibitem{MRW2008}
M.\,J.\ Mossinghoff, G.\ Rhin, Q.\ Wu,
\textit{Minimal Mahler measures},
Experiment.\ Math.\ \textbf{17} (2008), no.\ 4, 451--458.

\bibitem{RV1999}
F.\ Rodriguez Villegas,
\textit{Modular Mahler measures I},
in \textit{Topics in Number Theory (University Park, PA, 1997)},
Math.\ Appl.\ \textbf{467}, Kluwer (1999), 17--48.

\bibitem{RZ2014}
M.\ Rogers, W.\ Zudilin,
\textit{From $L$-series of elliptic curves to Mahler measures},
Compos.\ Math.\ \textbf{148} (2012), no.\ 2, 385--414.

\bibitem{Salem1944}
R.\ Salem,
\textit{A remarkable class of algebraic integers. Proof of a conjecture of Vijayaraghavan},
Duke Math.\ J.\ \textbf{11} (1944), 103--108.

\bibitem{Siegel1944}
C.\,L.\ Siegel,
\textit{Algebraic integers whose conjugates lie in the unit circle},
Duke Math.\ J.\ \textbf{11} (1944), 597--602.

\bibitem{Smyth1971}
C.\,J.\ Smyth,
\textit{On the product of the conjugates outside the unit circle of an algebraic integer},
Bull.\ London Math.\ Soc.\ \textbf{3} (1971), 169--175.

\bibitem{Smyth1981}
C.\,J.\ Smyth,
\textit{On measures of polynomials in several variables},
Bull.\ Austral.\ Math.\ Soc.\ \textbf{23} (1981), no.\ 1, 49--63.

\bibitem{Voutier1996}
P.\ Voutier,
\textit{An effective lower bound for the height of algebraic numbers},
Acta Arith.\ \textbf{74} (1996), no.\ 1, 81--95.

\end{thebibliography}

\end{document}
